% !TeX root = ../papers/nc-p3.tex
\section{有限な点対応の二窓の復元}
\label{sec:finite-point-chart}
四つの隣接点対応が有限な枝を扱う。系\ref{cor:finite-edge-point-sequences}により、
各$\Gamma_i$は長さ20で、共有する位置の射影はスキームとして同じ像を持つ。

\begin{proposition}[第三階数6の有限点列の四移動]\label{prop:rank-six-finite-four-shift}
標準AS対象で$\rho=6$、四つの隣接点対応が有限なら、四回の移動は点スキーム上の恒等射である。
\end{proposition}
\begin{proof}
隣接零点のどちらの端点でも、三テンソルの両端支持はともに3次元である
（命題\ref{prop:three-finite-edges}の証明）。中央支持は1でも2でもない。
隣接する二つの評価階数の和は6以下なので、全点列で各評価階数が3になる。
これは点列スキーム上でも定数階数である。実際、隣の評価行列の三次小行列式の単元で
階数3の直和因子を分裂させれば、積が零なので相手の四次小行列式はスキームとして消える。
従って二つの係数束の階数3の行空間は互いに零化空間である。
四移動した行空間は最初の行空間に戻り、その次の元の文字と再び隣接零点対をなす。
次の文字から前の文字への射が一意なので、元の文字に一致する。
この議論は直線束の局所自明化で行えるため非被約点の上でも成立する。
\end{proof}

\begin{corollary}[三辺からの周期点列]\label{cor:three-finite-periodic-sequences}
第三階数6で、連続する三辺$\Gamma_0,\Gamma_1,\Gamma_2$だけに有限性を仮定しても、
中央の長さ20のスキームの点列は両端へ一意に延び、この三辺の全零点を含む四周期点列を与える。
\end{corollary}
\begin{proof}
中央の辺で$b_1=f_2=3$となるのは前後の有限性だけから従う。
系\ref{cor:finite-edge-point-sequences}の局所直和因子の議論により、
中央の零点対$(x_1,x_2)$から前の$x_0$と次の$x_3$がスキーム上で一意に定まる。
両端の辺への射は共有文字への閉埋込みを通じて閉埋込みとなり、同じ長さ20なので同型である。
中央の二文字の両端支持は3,3で、評価もともに階数3となる。
従って$W_1,W_2$をその行空間とすれば$W_2=W_1^\perp$である。
端の二文字の行はそれぞれ$W_1^\perp,W_2^\perp$に入るから互いの積が零となり、
$(x_3,x_0)$は第四の辺の関係も満たす。ここでは端の評価の定数階数を必要としない。
四文字を繰り返せば結論を得る。第四の辺全体の有限性はこの証明からは従わない。
\end{proof}

\begin{theorem}[三つの座標空間からの復元と平方の直交]\label{thm:finite-square-chart}
標準AS対象で四辺の点対応が有限とし、移動で最初の三位置の点スキームを$X$に同一視する。
$D=H^0(X,\mathcal O_X)$とし、各座標直線束を自明化して
四次元の座標空間を$U_0,U_1,U_2\subset D$とする。
$UV$は積の線型包、$U^2=UU$と書く。次が成立する。
\begin{enumerate}
\item 隣接する$U,V$について$U^2,UV,V^2$は十次元で、それぞれの三乗$U^3,V^3$は$D$である。
$U^2 V^2$は超平面で、その零化汎関数$\tau$はスカラーを除き一意である。
$\tau(ab)$は完全対合であり、$U^2$と$V^2$は互いの直交補である。
\item 二つの対象が添字保持同型であることは、その$D$の代数同型と三つの独立な単元倍が
順序付き三空間を対応させることと同値である。
\item $U_2$を単元倍すると$U_2^2=U_0^2$とできる。従って第一窓を固定したとき、
次の座標空間は、$\mathbb P(D)$内の十本の二次式
$z^2\bmod U_0^2=0$が定める交叉に入る射影三平面のスキームから取れる。
\end{enumerate}
\end{theorem}
\begin{proof}
隣接対を$\mathbb P^3\times\mathbb P^3$の六本の双一次式の有限完全交叉と読む。
三次式の制限は定理\ref{thm:rho-four-positive-incidence}の証明の通り同型で、
一次・二次の制限も単射となる。$(1,1)$で消える切断は定義六式が張る。
実際、イデアルのKoszul分解を$(1,1)$でひねると、六本の定数切断のほかには
全次数0以下にコホモロジーを与える項がない。従って$\dim UV=16-6=10$であり、
$R_i$は対応する$U\otimes V\to D$の核そのものである。

$(2,2)$でひねった構造層のKoszul分解では、大域切断の列は次元$15,96,100$で、
ほかに外次数6の$H^6(\mathcal O(-4,-4))=k$が全次数0に出る。
他のコホモロジー項は零なので大域切断の制限像は余次元1である。
余核への辺写像は随伴同型$\omega_X=\mathcal O_X(2,2)$でのtraceである。
実際、Koszul分解を$\omega_{\mathbb P^3\times\mathbb P^3}$で双対化すると外次数$j$と$6-j$が対応し、
この写像はSerre双対の下で$k\to H^0(\mathcal O_X)$の単位に対応する。
従って自明化後に得る$\tau$の積対合は有限Gorenstein環$D$の完全対合である。
$\tau(U^2 V^2)=0$と次元から直交補の主張が従う。

添字保持同型は点列を対応させ、座標直線束の自明化の差は単元倍である。
逆に三空間が対応すれば二つの積の核も対応し、系\ref{cor:two-window-extension}から
AS延長の同型が従う。最初の三つの位置を順序付きで用いている。
最後に辺0と辺1のtraceを$\tau_0,\tau_1$とすると$\tau_1=\tau_0(g\,\cdot)$となる単元$g\in D^\times$がある。
従って$g U_2^2=U_0^2$である。各局所Artin因子では剰余体で平方根を取り、
冪零部分は有限二項展開で持ち上がるので$g$に平方根がある。
この単元で$U_2$を拡大縮小すればよい。偏極により三平面の条件は$U_2^2\subset U_0^2$と同値である。
候補から二窓を作った後には重なりと全次数AS条件を課す必要がある。
\end{proof}

\paragraph{孤立した第三の座標空間を持つ例。}
$v_0,\ldots,v_3$を一次式の基底、$v_4=\sum_{i=0}^3v_i$とし、四次元の二次式系
$\langle v_i^2-v_4^2:0\le i<4\rangle$を取る。共通零点があれば四つの符号の和が$\pm1$となるので、
この系は基点自由で定理\ref{thm:symmetric-frobenius-chart}の対称AS枝を与える。
辺の零点対は対称積がこの系に入る分解である。
$\sum_i t_i v_i^2$、$\sum_i t_i=0$で非零係数が三本または四本なら、その数が階数である。
五本なら$4\times5$の基底枠、可逆対角行列、その転置の積なので階数は3以上である。
従って階数2の二次式は十個の$v_i^2-v_j^2$だけで、
零点は$v_i+v_j,v_i-v_j$の順序を両方取る二十点である。完全交叉の長さから全て被約である。

ここでは$D=k^{20}$で、行をこの二十点とする評価行列から全積を成分ごとに計算できる。
$U_2=U_0$の平方を固定した三平面スキームの接空間は零である。
実際、四つの基底列$a$の変分$h(a)\in D$に対し、十組の
$a h(b)+b h(a)\bmod U_0^2$を零とする線型方程式は、80変数で階数64である。
残る16次元は基底変更で尽きる。
\path{computations/verify_finite_square_chart.py}で積・trace・接階数を有理数で検算できる。
従ってこの例で第一窓を固定すると、第三空間のこの正規化された解は被約孤立点である。
この例の他の解も含む全決定は定理\ref{thm:finite-all-square-roots}で行う。

\subsection{一つの乗法行列による平方根の探索}
\begin{proposition}[長さ20の環での比による復元]\label{prop:finite-square-ratio}
上の有限枝の$D$と十次元空間$Q=U_0^2$を固定する。
単元を含み、$L^2=Q$かつ$\operatorname{Sym}^3L\to D$が同型となる
四次元空間$L$を、次の方法で全て求められる。
$[\alpha]\in\mathbb P(D/k)$に対し
\[
 K_\alpha=\ker\bigl(Q\xrightarrow{\ \alpha\cdot\ }D/Q\bigr),\qquad
 J_\alpha=\{h\in D:hQ\subset K_\alpha^2\}
\]
を作る。十行十列の写像の階数が6、$\dim K_\alpha^2=10$で、
$J_\alpha$に単元$h$があれば$L=h^{-1/2}K_\alpha$を取り、上の単元と三乗の条件を検査する。
通過した場合$J_\alpha=kh$であり、一つの$[\alpha]$から得る空間は平方が1の単元倍を除いて一意である。
\end{proposition}
\begin{proof}
まず$L$を所与とし、単元$s\in L$と$s$に独立な$t\in L$を取り$\alpha=t/s$と置く。
$q\in K_\alpha$なら$q'=\alpha q\in Q$である。$S^2L\to Q$の同型で$q,q'$を二次式として読むと、
$sq'=tq$は$\operatorname{Sym}^3 L$での恒等式になる。
独立な二一次式は互いに素なので$q=s a$、$q'=t a$、$a\in L$である。
逆も成立し、$K_\alpha=sL$、$K_\alpha^2=s^2 Q$となる。

$gQ\subset Q$なら$g$は定数であることを示す。単元で$L$を割って$1\in L$としてよい。
$L=\langle 1,u,v,w\rangle$と書けば、$1,u,v,w$の斉次三次単項式の評価が基底であり、
同じことを$u,v,w$の次数3以下の単項式の独立性として読める。
$g\cdot1\in Q$より$g$は次数2以下で、$gu\in Q$より次数1以下、
$gu^2\in Q$より定数である。この議論の最初の二段は
$gL\subset L^2\Rightarrow g\in L$も与える（自明化を戻しても成立する）。

逆に上記の線型計算から単元$h$が得られれば、次元により$hQ=K_\alpha^2$で、
$h^{-1/2}K_\alpha$の平方は$Q$となる。
三乗等の条件を通過すれば$h^{-1}J_\alpha=k$で、単元の平方根の存在は前定理通りである。
さらに$s=h^{1/2}$とおくと$sL=K_\alpha\subset Q$かつ$\alpha sL\subset Q$なので、
先の含意から$s,\alpha s\in L$である。従って実際に二つの一次式の比を取っている。
$\alpha$への定数の加法と非零スカラー倍は$K_\alpha$を変えない。
一つの解に対応する比の族は、$L$の単元直線$ks$とそれを含む二平面$\langle s,t\rangle$の選択で
記述できる。二つの比の一致を交叉乗算して$\operatorname{Sym}^2L$の中で比較すれば
この旗は一意であり、そのパラメータは$3+2=5$次元である。
\end{proof}
この表示の探索変数は射影十八次元に入り、階数は小行列式で検査できる。
正規化した平方根を得た後に、二次関係・三次交叉・正則性を課す点は同じである。

有理数検算では、先の二十点例で比の乗法行列の階数6、$J_\alpha$の余次元19と正規化された平方を検査する。
さらに平方の復元だけの非被約な対照例として$D=k[t]/(t^{20})$で
$L=\langle\exp(at):a=0,1,4,16\rangle$を取る（指数関数は打ち切る）。
三重和二十個が相異なるので三乗はVandermonde行列によって$D$全体になる。
$\alpha=\exp(12t)$では乗法行列の核は$\exp(4t)L$である。
この対照例の検査は平方と比の計算に限る。

\begin{proposition}[比の階数条件に含まれる二つの成分]\label{prop:finite-ratio-components}
上の二十点の被約例で$Q=U_0^2$とする。
$\mathbb P(D/k)$内で比の乗法行列の階数が6以下となる軌跡には、
$U_0$の旗と$U_1$の旗から来る異なる二つの五次元既約成分がある。
後者の一般点では$J_\alpha=0$である。
\end{proposition}
\begin{proof}
前者の旗で階数が6であることは前命題による。後者については$Q$と$U_1^2$が
あるFrobenius積対合で互いに直交補であることを使う。
$Q\to D/Q$と$U_1^2\to D/U_1^2$の乗法行列はこの対合で互いに転置になるので、
$U_1$の比も同じ階数6を与える。いずれの旗の像も既約で五次元である。

各側で$s,t$の係数を順に$(1,3,9,27),(2,-1,4,7)$とする。
$\alpha=t/s$での階数条件の接方程式は
$\beta K_\alpha\subset Q+\alpha Q$である。実際、階数6の可逆小行列を分裂させると、
一次項は核から余核への写像だけである。通常の成分積と成分和による対合を使えば、
二十変数$\beta$への条件の階数は
$K_\alpha\,(Q+\alpha Q)^\perp$の次元であり、両側で13になる。
$k\cdot1+k\alpha$を除いた接次元は$20-13-2=5$である。
従って各試験点は滑らかで、その旗の像の閉包が既約成分となる。
双方で$\dim K_\alpha^2=10$で、$hQ\subset K_\alpha^2$の$h$に対する線型方程式は
前者で階数19、後者で20である。全て上記の検算ファイルで有理数計算できる。
$\dim J_\alpha\ge1$はこれらの定数階数開部分で閉条件であり、前者の旗で常に成立する。
後者の試験点では成立しないので二成分は異なる。
二成分以外の解の有無はこの接計算からは決まらない。
\end{proof}

この例では座標台から来る第三の五次元成分もある。
一般に上の真の比について、任意の$c\in k$で
$\dim D/(\alpha-c)D\le10$である。
実際、$t-cs$は非零一次式で、その$L^2$倍は三乗の独立性により十次元である。
$D$が被約なら二十座標のうち十一個以上が同じ値になる比を初めから除ける。

例の五つの一次式から辞書順の添字対を取り、各対で和と差、次に差と和を置いて
二十行を0から番号付けする。$S=\{0,2,3,7,8,16\}$上に支持される
$\alpha$の像は$\mathbb P(D/k)$の射影五平面で、乗法の階数は高々6である。
$\alpha$を$S$の指示関数とすると階数は6、接方程式の階数は13であり、
この五平面は再び既約成分である。これは台の補で同じ値を持つので
$U_0,U_1$いずれの旗からも稠密には覆えない。この成分の除外は単元の検査でも行える。
